\documentclass[12pt,a4paper,openany]{report}

% =========================================================
% PACKAGES
% =========================================================
\usepackage[a4paper,margin=1in]{geometry}
\usepackage{setspace}
\usepackage{graphicx}
\usepackage{float}
\usepackage{booktabs}
\usepackage{array}
\usepackage{longtable}
\usepackage{xcolor}
\usepackage{listings}
\usepackage{hyperref}
\usepackage{enumitem}
\usepackage{fancyhdr}
\usepackage{amsmath}
\usepackage{tikz}

% =========================================================
% GENERAL SETTINGS
% =========================================================
\onehalfspacing

\hypersetup{
    colorlinks=true,
    linkcolor=black,
    urlcolor=blue,
    citecolor=black
}

% =========================================================
% CODE LISTING STYLE
% =========================================================
\lstset{
    language=Python,
    basicstyle=\ttfamily\small,
    keywordstyle=\color{blue},
    commentstyle=\color{gray},
    stringstyle=\color{red},
    numbers=left,
    numberstyle=\tiny,
    stepnumber=1,
    numbersep=8pt,
    frame=single,
    breaklines=true,
    showstringspaces=false,
    tabsize=4
}

% =========================================================
% HEADER AND FOOTER
% =========================================================
\pagestyle{fancy}
\fancyhf{}
\fancyhead[L]{Software Maintenance Lab Project}
\fancyhead[R]{Math Problem Solver}
\fancyfoot[C]{\thepage}

% =========================================================
% DOCUMENT
% =========================================================
\begin{document}

% =========================================================
% TITLE PAGE
% =========================================================

\begin{titlepage}
    \centering

    \vspace*{1.0cm}

    {\Large \textbf{ISLAMIC UNIVERSITY OF TECHNOLOGY}}\\[0.4cm]
    {\large \textbf{Department of Computer Science and Engineering}}\\[1.5cm]

    {\Huge \textbf{Software Maintenance Lab Project Report}}\\[1.0cm]

    {\LARGE \textbf{Math Problem Solver}}\\[1.8cm]

    \renewcommand{\arraystretch}{1.5}

    \begin{tabular}{ll}
        \textbf{Course Code:} & SWE-4802 \\
        \textbf{Course Title:} & Software Maintenance Lab \\
        \textbf{Name:} & Muhammad Yasir Raiyan \\
        \textbf{ID:} & 210042152 \\
        \textbf{Dept:} & CSE \\
        \textbf{Prog:} & B.Sc. in S.W.E \\
        \textbf{Sem:} & 8th \\
        \textbf{Academic Year:} & 2024--25 \\
    \end{tabular}

    \vfill

    {\large \textbf{Software Maintenance Project}}

    \vspace{0.5cm}

    {\large 2026}

\end{titlepage}

% =========================================================
% TABLE OF CONTENTS
% =========================================================

\pagenumbering{roman}

\tableofcontents

% IMPORTANT:
% No \clearpage here.
% This prevents an unnecessary blank page.

\pagenumbering{arabic}

% =========================================================
% CHAPTER 1
% =========================================================

\chapter{Project Overview}

\section{Introduction}

Math Problem Solver is a Python-based mathematical calculation
application. The initial version of the application provides basic and
advanced mathematical operations through a command-line interface.
The system was subsequently adapted and improved through different
software maintenance activities.

The main objective of this project is to simulate the four major
categories of software maintenance:

\begin{enumerate}
    \item Corrective Maintenance
    \item Adaptive Maintenance
    \item Preventive Maintenance
    \item Perfective Maintenance
\end{enumerate}

For each maintenance category, the following activities were performed:

\begin{enumerate}
    \item Program Comprehension
    \item Change Management
    \item Impact Analysis
    \item Reverse Engineering
    \item Refactoring
\end{enumerate}

The project demonstrates how an existing software system can be
understood, modified, improved, tested, and evolved throughout its
maintenance lifecycle.

\section{Project Objectives}

The major objectives of this project are:

\begin{itemize}
    \item Identify and correct existing software defects.
    \item Adapt the existing system to a new execution environment.
    \item Improve maintainability and prevent future defects.
    \item Add new functionality based on user requirements.
    \item Analyze the impact of software changes.
    \item Perform reverse engineering of existing functionality.
    \item Refactor existing code without changing intended behavior.
    \item Maintain regression testing throughout the maintenance
          lifecycle.
\end{itemize}

% =========================================================
% CHAPTER 2
% =========================================================

\chapter{Initial System -- V1 Baseline}

\section{System Description}

The initial version of the Math Problem Solver provides mathematical
calculation functionality through a command-line interface.

The baseline system contains:

\begin{itemize}
    \item Basic arithmetic operations
    \item Advanced mathematical operations
    \item Input validation
    \item Result formatting
    \item Logging
    \item Unit testing
    \item Integration testing
\end{itemize}

\section{Initial Project Structure}

The initial project structure is shown below:

\begin{lstlisting}[language={}]
Math-Problem-Solver/
|
+-- app.py
+-- main.py
|
+-- calculator/
|   +-- basic.py
|   +-- advanced.py
|   +-- validator.py
|
+-- utils/
|   +-- logger.py
|   +-- formatter.py
|
+-- tests/
|   +-- test_basic.py
|   +-- test_advanced.py
|   +-- test_validator.py
|   +-- test_integration.py
|
+-- templates/
|   +-- index.html
|
+-- requirements.txt
+-- README.md
\end{lstlisting}

\section{Baseline Testing}

Before performing maintenance activities, the initial version was
tested.

\begin{verbatim}
Ran 10 tests
OK
\end{verbatim}

The successful baseline test result was used as the reference point
for subsequent maintenance activities.

% =========================================================
% CHAPTER 3
% =========================================================

\chapter{Corrective Maintenance}

\section{Scenario}

A defect was identified in the division operation. When the denominator
was zero, the application produced a \texttt{ZeroDivisionError}, which
caused the application to terminate unexpectedly.

\section{Change ID}

\textbf{CM-001}

\section{Program Comprehension}

The existing calculation flow was analyzed to locate the source of the
problem.

\begin{verbatim}
User
  |
  v
main.py
  |
  v
calculator/basic.py
  |
  v
divide(a, b)
  |
  v
a / b
\end{verbatim}

The \texttt{divide()} function was identified as the source of the
division-by-zero defect.

\section{Change Management}

\subsection{Problem}

Division by zero caused the application to crash with a
\texttt{ZeroDivisionError}.

\subsection{Requested Change}

The system should prevent division by zero and provide a meaningful
error message to the user.

\subsection{Change Implemented}

The \texttt{divide()} function was modified to validate the denominator
before performing the calculation.

\section{Impact Analysis}

The following components were affected:

\begin{table}[H]
\centering
\renewcommand{\arraystretch}{1.3}
\begin{tabular}{p{5cm} p{3cm}}
\toprule
\textbf{Component} & \textbf{Impact} \\
\midrule
\texttt{calculator/basic.py} & Direct Impact \\
\texttt{main.py} & Indirect Impact \\
\texttt{tests/test\_basic.py} & Regression Test \\
\bottomrule
\end{tabular}
\caption{Corrective Maintenance Impact Analysis}
\end{table}

\section{Reverse Engineering}

The original execution flow was reconstructed as follows:

\begin{verbatim}
User
  |
  v
Select Division
  |
  v
main.py
  |
  v
divide(a, b)
  |
  v
Denominator = 0
  |
  v
ZeroDivisionError
\end{verbatim}

This analysis helped identify where the defect occurred and which
components required modification.

\section{Refactoring}

The division function was changed to:

\begin{lstlisting}
def divide(a: float, b: float) -> float:
    if b == 0:
        raise ValueError("Cannot divide by zero")
    return a / b
\end{lstlisting}

A regression test was added:

\begin{lstlisting}
def test_divide_by_zero(self):
    with self.assertRaises(ValueError):
        divide(10, 0)
\end{lstlisting}

\section{Testing Result}

After implementing the corrective maintenance:

\begin{verbatim}
Ran 11 tests
OK
\end{verbatim}

\section{Version Transition}

\begin{center}
\textbf{V1}
$\rightarrow$
\textbf{Corrective Maintenance}
$\rightarrow$
\textbf{V1.1}
\end{center}

% =========================================================
% CHAPTER 4
% =========================================================

\chapter{Adaptive Maintenance}

\section{Scenario}

The original application was command-line based. The system needed to
be adapted so that users could access the calculator through a web
browser.

\section{Change ID}

\textbf{CM-002}

\section{Program Comprehension}

The existing architecture was analyzed:

\begin{verbatim}
User
  |
  v
main.py
  |
  v
Calculator Modules
  |
  v
Result
\end{verbatim}

The calculation modules were already separated from the main
application logic. Therefore, they could be reused when adapting the
system to a web environment.

\section{Change Management}

\subsection{Requirement}

The existing command-line application needed to be adapted to a
web-based environment.

\subsection{Change Implemented}

A Flask web interface was introduced while reusing the existing
calculator modules.

\subsection{Technology Added}

\begin{itemize}
    \item Flask
    \item HTML
    \item Web interface
\end{itemize}

\section{Impact Analysis}

The affected components were:

\begin{table}[H]
\centering
\renewcommand{\arraystretch}{1.3}
\begin{tabular}{p{5cm} p{3cm}}
\toprule
\textbf{Component} & \textbf{Impact} \\
\midrule
\texttt{app.py} & Web Application \\
\texttt{requirements.txt} & Flask Dependency \\
\texttt{templates/index.html} & Web Interface \\
\bottomrule
\end{tabular}
\caption{Adaptive Maintenance Impact Analysis}
\end{table}

The following components were reused:

\begin{itemize}
    \item \texttt{calculator/basic.py}
    \item \texttt{calculator/advanced.py}
    \item \texttt{calculator/validator.py}
\end{itemize}

\section{Reverse Engineering}

The original flow was:

\begin{verbatim}
CLI
 |
 v
main.py
 |
 v
Validation
 |
 v
Calculator Function
 |
 v
Result
\end{verbatim}

The adapted flow became:

\begin{verbatim}
Browser
 |
 v
index.html
 |
 v
app.py
 |
 v
Validation
 |
 v
Calculator Function
 |
 v
Result
 |
 v
Browser
\end{verbatim}

\section{Refactoring}

A Flask web interface was added without rewriting the existing
calculator logic.

The existing calculator functions were reused by \texttt{app.py}.

\begin{verbatim}
Browser
  |
  v
Flask app.py
  |
  v
calculator/
  |-- basic.py
  |-- advanced.py
  +-- validator.py
  |
  v
Result
\end{verbatim}

The Math Problem Solver can now be accessed through a web browser.

\section{Version Transition}

\begin{center}
\textbf{V1.1}
$\rightarrow$
\textbf{Adaptive Maintenance}
$\rightarrow$
\textbf{V1.2}
\end{center}

% =========================================================
% CHAPTER 5
% =========================================================

\chapter{Preventive Maintenance}

\section{Scenario}

The application was working correctly, but improvements were required
to prevent future defects and make the system easier to maintain.

\section{Change ID}

\textbf{CM-003}

\section{Program Comprehension}

The existing project structure and dependencies were analyzed.

The following areas were identified for improvement:

\begin{itemize}
    \item Input validation
    \item Error handling
    \item Code readability
    \item Type information
    \item Test coverage
\end{itemize}

\section{Change Management}

\subsection{Objective}

The objective was to improve the maintainability, reliability, and
readability of the system.

\subsection{Changes Implemented}

The following improvements were performed:

\begin{itemize}
    \item Improved input validation
    \item Added type hints
    \item Improved error handling
    \item Added mathematical input validation
    \item Added additional tests
\end{itemize}

\section{Impact Analysis}

The affected components were:

\begin{table}[H]
\centering
\renewcommand{\arraystretch}{1.3}
\begin{tabular}{p{6cm} p{3cm}}
\toprule
\textbf{Component} & \textbf{Impact} \\
\midrule
\texttt{calculator/validator.py} & High \\
\texttt{calculator/basic.py} & Medium \\
\texttt{calculator/advanced.py} & Medium \\
\texttt{tests/test\_validator.py} & High \\
\texttt{tests/test\_advanced.py} & High \\
\bottomrule
\end{tabular}
\caption{Preventive Maintenance Impact Analysis}
\end{table}

The intended functionality of the calculator was preserved.

\section{Reverse Engineering}

The validation flow was analyzed:

\begin{verbatim}
User Input
 |
 v
validate_number()
 |
 v
Convert to number
 |
 v
Calculator
 |
 v
Result
\end{verbatim}

Potential invalid inputs were identified before they could cause
unexpected behavior.

\section{Refactoring}

The validation function was improved with type hints and better
exception handling:

\begin{lstlisting}
def validate_number(value: str) -> float:
    try:
        return float(value)
    except (ValueError, TypeError):
        raise ValueError(
            "Invalid number. Please enter a valid number."
        )
\end{lstlisting}

The advanced calculator was also improved to handle invalid square-root
and factorial inputs.

Additional validation tests were added.

\section{Result}

The software became easier to understand, maintain, and extend while
preserving its existing behavior.

\section{Version Transition}

\begin{center}
\textbf{V1.2}
$\rightarrow$
\textbf{Preventive Maintenance}
$\rightarrow$
\textbf{V1.3}
\end{center}

% =========================================================
% CHAPTER 6
% =========================================================

\chapter{Perfective Maintenance}

\section{Scenario}

Users requested additional mathematical functionality and an improved
web experience.

\section{Change ID}

\textbf{CM-004}

\section{Program Comprehension}

The existing advanced calculator and web application were analyzed to
determine where the new features should be integrated.

\begin{verbatim}
Browser
 |
 v
app.py
 |
 v
Calculator Module
 |
 v
Result
\end{verbatim}

\section{Change Management}

\subsection{User Requirements}

Users requested:

\begin{itemize}
    \item Percentage calculation
    \item Factorial calculation
    \item Improved operation selection
\end{itemize}

\subsection{Objective}

The objective was to improve the functionality and user experience of
the application.

\section{Impact Analysis}

The following components were affected:

\begin{itemize}
    \item \texttt{calculator/advanced.py}
    \item \texttt{app.py}
    \item \texttt{templates/index.html}
    \item \texttt{tests/test\_advanced.py}
\end{itemize}

Existing basic calculator functionality remained unchanged.

\section{Reverse Engineering}

The existing feature flow was analyzed:

\begin{verbatim}
Browser
 |
 v
app.py
 |
 v
Operation Selection
 |
 v
Calculator Function
 |
 v
Result
\end{verbatim}

The new features were integrated into the existing architecture.

\section{Refactoring}

Two new mathematical operations were added.

\subsection{Factorial}

\begin{lstlisting}
def factorial(n: int) -> int:
    if n < 0:
        raise ValueError(
            "Factorial is not defined for negative numbers."
        )

    if not float(n).is_integer():
        raise ValueError(
            "Factorial requires a whole number."
        )

    return math.factorial(int(n))
\end{lstlisting}

\subsection{Percentage}

\begin{lstlisting}
def percentage(value: float, percent: float) -> float:
    return (value * percent) / 100
\end{lstlisting}

The web interface was updated to allow users to select the new
operations.

\section{Result}

The final application supports:

\begin{itemize}
    \item Addition
    \item Subtraction
    \item Multiplication
    \item Division
    \item Power
    \item Square Root
    \item Cube Root
    \item Factorial
    \item Percentage
\end{itemize}

\section{Testing Result}

After perfective maintenance:

\begin{verbatim}
Ran 16 tests
OK
\end{verbatim}

\section{Version Transition}

\begin{center}
\textbf{V1.3}
$\rightarrow$
\textbf{Perfective Maintenance}
$\rightarrow$
\textbf{V2.0}
\end{center}

% =========================================================
% CHAPTER 7
% =========================================================

\chapter{Overall Maintenance History}

The complete maintenance lifecycle is shown below:

\begin{verbatim}
Initial V1 Baseline
        |
        v
Corrective Maintenance
Fix division by zero
        |
        v
V1.1
        |
        v
Adaptive Maintenance
Add Flask web interface
        |
        v
V1.2
        |
        v
Preventive Maintenance
Improve validation and code quality
        |
        v
V1.3
        |
        v
Perfective Maintenance
Add percentage and factorial
        |
        v
V2.0
\end{verbatim}

The project evolved incrementally from the initial baseline to the
final maintained version.

% =========================================================
% CHAPTER 8
% =========================================================

\chapter{Testing Results}

Testing was performed throughout the software maintenance lifecycle to
ensure that modifications did not break existing functionality.

\section{Testing Summary}

\begin{table}[H]
\centering
\renewcommand{\arraystretch}{1.4}
\begin{tabular}{p{6cm} p{4cm} p{2cm}}
\toprule
\textbf{Stage} & \textbf{Result} & \textbf{Status} \\
\midrule
Initial V1 Baseline & 10 tests passed & PASS \\
Corrective Maintenance & 11 tests passed & PASS \\
Preventive Maintenance & Tests passed & PASS \\
Perfective Maintenance & 16 tests passed & PASS \\
\bottomrule
\end{tabular}
\caption{Testing Results}
\end{table}

\section{Final Test Command}

The final test suite was executed using:

\begin{lstlisting}[language={}]
python -m unittest discover -s tests -p "test_*.py"
\end{lstlisting}

\section{Final Test Result}

\begin{verbatim}
Ran 16 tests
OK
\end{verbatim}

The successful execution of all 16 tests demonstrates that the final
system maintained its expected functionality after the maintenance
activities.

% =========================================================
% CHAPTER 9
% =========================================================

\chapter{Git Maintenance History}

Git was used to organize the software maintenance activities into
separate logical commits.

The maintenance history was organized as follows:

\begin{verbatim}
Initial V1 baseline
        |
        v
Corrective maintenance:
fix division by zero
        |
        v
Adaptive maintenance:
add Flask web interface
        |
        v
Preventive maintenance:
improve validation and code quality
        |
        v
Perfective maintenance:
add percentage and factorial
        |
        v
Complete software maintenance documentation
\end{verbatim}

This commit-based organization makes it easier to identify, track, and
review each maintenance activity independently.

% =========================================================
% CHAPTER 10
% =========================================================

\chapter{Comparison of Maintenance Categories}

\begin{table}[H]
\centering
\renewcommand{\arraystretch}{1.5}
\begin{tabular}{p{3.2cm} p{4cm} p{5.2cm}}
\toprule
\textbf{Category} &
\textbf{Purpose} &
\textbf{Project Change} \\
\midrule

Corrective &
Fix existing defects &
Fixed division by zero \\

Adaptive &
Adapt to a new environment &
Added Flask web interface \\

Preventive &
Prevent future problems &
Improved validation, type hints, error handling, and testing \\

Perfective &
Improve functionality &
Added factorial and percentage \\

\bottomrule
\end{tabular}
\caption{Comparison of Maintenance Categories}
\end{table}

% =========================================================
% CHAPTER 11
% =========================================================

\chapter{Conclusion}

The Math Problem Solver was successfully used to simulate the four major
categories of software maintenance.

Corrective maintenance fixed an existing division-by-zero defect.
Adaptive maintenance adapted the application from a command-line
environment to a web-based environment.

Preventive maintenance improved validation, error handling, type hints,
and testing. Finally, perfective maintenance enhanced the application
by adding factorial and percentage calculations.

The maintenance activities were performed systematically using program
comprehension, change management, impact analysis, reverse engineering,
and refactoring.

The final system successfully passed:

\begin{center}
\Large
\textbf{16 Tests -- OK}
\end{center}

Therefore, the project demonstrates a complete software maintenance
lifecycle covering corrective, adaptive, preventive, and perfective
maintenance activities.

% =========================================================
% END DOCUMENT
% =========================================================

\end{document}
